%!TEX root = ../rapport.tex

\chapter{Future Work}

We have implemented but a small part of what one could imagine a program like
this to do. Much more work can be done; for instance, the grapical user can
surely be greatly enhanced interface, since our focus has not been on
designing a user interface. Moreover, features relevant for the actual
drawings can also be implemented, not to mention performance tweaks of
different kinds. In this Chapter we give an (inherently incomplete) list of
what remains to be done.

\begin{description}
	\item[Multithreading.] Except for the graphics widgets, with whose internal workings 
	we have
	nothing to do, the current implementation works in a single thread. That means
	that the drawing routines have to wait for the graph generator to finish the
	current graph, which is a waiting period that can be felt for big graphs, $V>100$ e.g.
	In a future version, there is a worker thread that emits graphs to a list, making 
	it ``look ahead'' such that when the user pushes the ``Forward''-button the graph
	is already ready for the layout algorithm to work on. This would greatly increase
	the overall smoothness of the program. 
	
	\item[Zooming.] An crucial feature for very large, cluttered graphs is the
	ability to zoom in on selected areas and pan the viewing window. This could be combinated
	with an ability to intelligently switch drawing algorithms for certain areas. 
	If a slow, but good-looking, algorithm is at hand, for instance, it should be used
	on the small subgraph that the user has zoomed in on. 
	
	\item[3D.] Some of the algorithms (Neato and Majorization) currently implemented generalizes to 
	$n$ dimensions. Thus, it would be interesting to see how three dimensional 
	graph embeddings look like, and whether they can be beneficial to use in some
	circumstances. In the current implementation, the Cairo library for two-dimensional
	vector drawing is used, so in a version with three-dimensional drawings another
	graphics library that supports 3D must be used. Since the drawing routines 
	only use the graphics primitives served by Cairo, it would not be a big challenge
	to switch the routines to use the primitives from e.g. OpenGL instead \cite{OpenGL}.
	
	\item[Exporting facilities.] The ability to easily take screenshots through 
	a push of button, for instance, would be quite useful to have. Furthermore,
	if one were able to export the drawn graph to other formats than just image files,
	an integration with existing graph software could be established. Several
	attempts at standard graph description languages/formats have been suggested
	\cite{GraphML,GXL,DOT,XGMML}, and since none of them have gained universal 
	acceptance as ``the'' standard, export facilities for all or most of them 
	should be implemented .
	
	\item[Drawing algorithms.] More drawing algorithms should be implemented. 
	The algorithm by Civril et al., for instance, is claimed to be efficient for
	large graphs (for $V>20000$ it uses approximately 100 seconds) \cite{Civril2005}.
	
	\item[Memory efficiency.] More work can be done to optimize memory usage. 
	For instance, when stepping through several generations of graphs, all of
	them are kept in memory until a new push of the ``Draw''-button. A possibility
	to set a limit of the amount of graphs remembered in memory could help in reducing
	the memory footprint, at the expense of a risk of recomputing some of the graphs
	if the user steps backwards.
	
	\item[GUI.] Much more can be done to enhance the graphical user interface. 
	As of now, the program does not contain a menu bar for instance. If new features
	like a properties window and exporting facilities are added, they should be
	placed in a menu entry, since they are not directly relevant for the actual
	graph generation and -drawing. 
	
	In said properties window, options to control memory usage, graph colour codes,
	parameters for the random lambda term generator etc. should be placed.
	
	\item[Better redex highlighting.] When a node representing a lambda term with 
	many redexes is clicked, all of its outgoing edges are highlighted in the same
	way. It would be helpful if there were some kind of differentiation in the way
	these are highlighted, such that the user easily can see which edge represent
	which redex. We imagine that this could be done in a way where the edges and their
	corresponding redexes in the textual representation of the term are given
	matching colours, where each edge/redex-pair is coloured with its own, distinct
	colour.
	
\end{description}

% \begin{itemize}
	% \item Multithreaded
	% % \item Dragging of nodes
	% % \item Hovering/clicking actions
	% \item Zooming/panning
	% \item 3D-view
	% \item Export of graph as image file, as dot file and as pickle Python object
	% \item Calculation of the next generations of graphs in a background thread
	% \item Implement more layout algorithms (e.g. confluent drawings).
	% \item Clever handling of memory usage. If too much RAM is used, be able
	% to free the ``oldest'' graphs in the list to make more room.
	% \item Different colouring of highlighted edges with the same colouring of
	% the redexes in the term.
	% \item Menubar, preferences, screenshot etc.
	% ???	
% \end{itemize}